\section{Methodology}
% [WRITING GUIDANCE]
% Target length: 1.5 - 2 columns.
% Discuss the CRISP-DM phases.

\subsection{Data Preprocessing and Feature Engineering}
% Summarize the preprocessing steps (imputation, outlier removal).
% Discuss the creation of composite business metrics (Customer_Value_Score, etc.).

% Insert Feature Engineering Summary Table
\begin{table}[htbp]
\caption{Feature Engineering Summary}
\begin{center}
\begin{tabular}{|l|p{5cm}|}
\hline
\textbf{New Feature} & \textbf{Derivation Logic} \\
\hline
Total\_Spending & Sum of all product category spending \\
Customer\_Value\_Score & Formula combining spending, acceptance, and recency \\
\hline
\end{tabular}
\label{tab:feature_eng}
\end{center}
\end{table}

\subsection{Feature Selection}
% Discuss methods used to reduce dimensionality (collinearity checks, RFE).

\subsection{Classification Algorithms}
% Describe the models compared: Decision Tree, Naive Bayes (GaussianNB), k-NN, Random Forest, XGBoost, etc.
% For Naive Bayes, state its conditional independence assumption.
% For k-NN, mention the distance metric and scaling requirement.

\subsection{Evaluation Metrics}
% Discuss why Accuracy is insufficient due to imbalance.
% Introduce F1-Score, ROC-AUC, and PR-AUC.

% Write your methodology details here...
